\chapter{Robustness and Sensitivity Analysis}
\label{ch:robustness}

\section{Lag Sensitivity}

The choice of lag order in the ADL model could affect the results. To verify that the main findings are not an artefact of the selected lag structure, I estimate the model over a grid of all combinations of $p \in \{1, 2, 3, 4\}$ and $q \in \{0, 1, 2, 3\}$, giving 16 specifications. All models are estimated on the same sample for comparability.

Two patterns emerge from the grid. First, the positive cumulative pass-through (CPT$^+$) is near zero when no oil lags are included ($q = 0$) and increases as more lags are added, reaching about 0.019--0.021 at $q = 3$ regardless of the AR lag order. This stability across different AR orders is reassuring. Second, the asymmetry $p$-value ranges from 0.22 to 0.31 across the $q = 3$ specifications, confirming that the failure to reject symmetry is not sensitive to the number of CPI lags.

The primary model (ADL(3,3)) has the lowest AIC among the $q = 3$ specifications and ranks 13th overall. The specifications with the lowest AIC values tend to be those with $q = 0$, which estimate symmetric effects only. This makes sense: if asymmetry is not statistically strong, adding the extra positive and negative oil lags costs degrees of freedom without proportionately improving fit.

\section{Brent Plus Exchange Rate Specification}

Table~\ref{tab:brentexr} compares the primary model with the specification that separates Brent crude prices (in USD) from the rupee--dollar exchange rate.

\begin{table}[H]
\centering
\caption{Comparison: Primary INR Oil vs.\ Brent + Exchange Rate Specification}
\label{tab:brentexr}
\begin{tabular}{lcc}
\toprule
& \textbf{Primary (Oil\_INR)} & \textbf{Brent USD + EXR} \\
\midrule
CPT$^+$ & 0.0213 & 0.0275 \\
$p$(CPT$^+ = 0$) & 0.122 & 0.093 \\
CPT$^-$ & 0.0006 & $-0.0021$ \\
$p$(CPT$^- = 0$) & 0.938 & 0.779 \\
Asymmetry $p$ & 0.241 & 0.144 \\
EXR coefficient & -- & 0.0396 \\
$p$(EXR) & -- & 0.029 \\
Adj.\ $R^2$ & 0.449 & 0.458 \\
AIC & 440.77 & 438.79 \\
\bottomrule
\end{tabular}
\begin{flushleft}
\vspace{0.3em}
\small\textit{Notes:} Both models include identical AR lags, IIP, policy dummies, and monthly dummies.
\end{flushleft}
\end{table}

The Brent-plus-EXR specification produces a higher positive pass-through estimate (CPT$^+ = 0.028$), which is marginally significant at 10 per cent ($p = 0.093$). The exchange rate coefficient is positive and significant at 5 per cent ($p = 0.029$), confirming that rupee depreciation independently raises CPI inflation. A 1 per cent depreciation of the rupee raises monthly CPI inflation by about 0.04 percentage points, holding oil prices constant.

The asymmetry $p$-value drops from 0.241 to 0.144 in the Brent-plus-EXR specification. While still above the conventional 5 per cent threshold, this improvement suggests that separating the oil and exchange rate channels sharpens the asymmetry signal. The Brent-plus-EXR model also has a lower AIC (438.79 vs.\ 440.77), indicating marginally better overall fit.

\section{COVID Sensitivity}

The COVID-19 lockdown in April 2020 produced an extreme negative oil price shock (Brent fell about 45 per cent in a single month) alongside a sharp contraction in domestic activity. To verify that this single observation does not drive the results, I re-estimate the main model without the COVID dummy.

\begin{table}[H]
\centering
\caption{COVID Sensitivity Check}
\label{tab:covid}
\begin{tabular}{lccc}
\toprule
\textbf{Specification} & \textbf{CPT$^+$} & \textbf{CPT$^-$} & \textbf{Asym.\ $p$} \\
\midrule
With COVID dummy & 0.0213 & 0.0006 & 0.241 \\
Without COVID dummy & 0.0206 & 0.0017 & 0.267 \\
\bottomrule
\end{tabular}
\end{table}

The results are nearly identical. CPT$^+$ changes by less than 0.001, and the asymmetry $p$-value moves from 0.241 to 0.267. The COVID outlier does not materially affect the conclusions.

\section{Winsorized Oil Shocks}

As a further check against the influence of extreme oil price movements, I winsorise the positive and negative oil shock components at the 1st and 99th percentiles. This trims the most extreme months of oil price change and re-estimates the model with the trimmed values.

\begin{table}[H]
\centering
\caption{Winsorised Oil Shocks (1\% Tails)}
\label{tab:winsor}
\begin{tabular}{lccc}
\toprule
\textbf{Specification} & \textbf{CPT$^+$} & \textbf{CPT$^-$} & \textbf{Asym.\ $p$} \\
\midrule
Original & 0.0213 & 0.0006 & 0.241 \\
Winsorised (1\%) & 0.0226 & $-0.0004$ & 0.231 \\
\bottomrule
\end{tabular}
\end{table}

Winsorisation slightly increases the positive pass-through and pushes the negative pass-through marginally below zero, which is the expected direction. The asymmetry $p$-value remains essentially unchanged at 0.231. The main finding is robust to extreme oil price observations.

\section{Rolling Window Analysis}

To assess whether the oil pass-through relationship has changed over time, I estimate the main asymmetric model on a rolling 60-month window and track the cumulative positive and negative pass-through estimates across the sample.

\begin{figure}[H]
\centering
\includegraphics[width=0.9\textwidth]{fig_7_rolling_window.png}
\caption{}
\label{fig:rolling}
\end{figure}

Figure~\ref{fig:rolling} shows that the positive pass-through estimate varies over time but remains predominantly positive. It is highest during windows that include the 2007--08 oil price spike and lowest during the relatively calm period of 2015--17. The negative pass-through estimate fluctuates around zero throughout, consistent with the full-sample finding that oil price decreases have little effect on headline CPI.

\section{Residual Diagnostics}

Figure~\ref{fig:residuals} provides visual diagnostic checks for the main model's residuals.

\begin{figure}[H]
\centering
\includegraphics[width=0.95\textwidth]{fig_8_residual_diagnostics.png}
\caption{}
\label{fig:residuals}
\end{figure}

\section{Oil Price Regimes}

Figure~\ref{fig:regimes} plots the Brent crude oil price with shaded bands indicating distinct price regimes over the sample period. The visual makes clear that the sample covers substantial variation in oil price levels and volatility, which is necessary for estimating pass-through effects reliably.

\begin{figure}[H]
\centering
\includegraphics[width=0.9\textwidth]{fig_9_oil_price_regimes.png}
\caption{}
\label{fig:regimes}
\end{figure}

\begin{figure}[H]
\centering
\includegraphics[width=0.8\textwidth]{fig_10_asymmetry_gap.png}
\caption{}
\label{fig:asymgap}
\end{figure}

Figure~\ref{fig:asymgap} visually compares the positive and absolute negative pass-through estimates across the primary and robustness specifications. In every specification, CPT$^+$ exceeds $|$CPT$^-|$, reinforcing the conclusion that the asymmetry is a consistent feature of the point estimates, even if it falls short of statistical significance.
